\clearpage
\hypertarget{native-446-part-lower}{}
\begin{center}
{\large MONDSBEOBACHTUNGEN.}

\noindent\rule{\linewidth}{0.4pt}

Astronomische Nachrichten. Band I. Nr. 12. S. 189. April 1822.

\noindent\rule{\linewidth}{0.4pt}
\end{center}

Von Hofrat und Ritter \textsc{Gauss} habe ich die folgenden an der Göttinger Sternwarte beobachteten Unterschiede in der Rektaszension zwischen dem Mond und den bestimmten Sternen erhalten.

\begin{center}
\begin{tabular}{l|l|l|l}
1822 März 1 & Mond 1 R & & 7 Fäden \\
 & PV, 287 & $+9^{\mathrm{m}}\;54^{\mathrm{s}}73$ & 7 \\
 & V, 303 & $+12\quad55.11$ & 7 \\
 & V, 319 & $+15\quad23.77$ & 7 \\
März 3 & 309 Mayer & $-13\quad3.60$ & 7 \\
 & Mond & & 7 \\
 & $\mu$ Cancri & $+9\quad55.94$ & 6 \\
März 4 & 351 Mayer & $-15\quad25.29$ & 7 \\
 & 44 Cancri & $-9\quad58.83$ & 7 \\
 & $\delta$ Cancri & $-8\quad25.67$ & 7 \\
 & 371 Mayer & $-2\quad24\;57$ & 7 \\
 & Mond & & 7 \\
\end{tabular}
\end{center}

Zu den Beobachtungen von Kapitän v. \textsc{Caroc} in Kopenhagen an der Sternwarte der Gradmessung auf Holkens Bastion ist noch Folgendes hinzuzufügen.

\begin{center}
\begin{tabular}{l|l|l|l}
März 9 & Mond 2 R & & 3 Fäden \\
 & 56 Virginis & $+20^{\mathrm{m}}\;17^{\mathrm{s}}91$ & 5 \\
 & XIII, 56 & $+25\quad59.83$ & 5 \\
\end{tabular}
\end{center}

\begin{flushright}
\textsc{Schumacher}.
\end{flushright}
